\begin{frame}[plain]
    \begin{center}
        \LARGE Arquiteturas de redes neurais convolucionais    
    \end{center}
\end{frame}

\begin{frame}{\textit{ResNet152V2}}\label{sec:ResNet152V2}
    A arquitetura foi proposta por \textcite{he2016identitymappingsdeepresidual} e possui mais de $100$ camadas, a principal ideia foi utilizar função residual aditivo $\mathcal{F}$ em relação ao $h(x_{l})$ e com uma identidade $h(x_{l})=x_{l}$.
\end{frame}

\begin{frame}{\textit{ResNet152V2}}
    Arquitetura possui unidades denominados \textit{residual units} que podem ser expressados dessa forma:
\begin{align}\label{eq:residual-units}
    y_1=&h(x_l)+F(x_1,W_1)\\
    x_{l+1}=&f(y_1)
\end{align}

    Onde:
    \begin{itemize}
        \item $x_l$ e $x_{l+1}$ são as entradas e saídas da unidade de camada $l$
        \item $F$ é a função residual 
        \item $f$ é a ativação \textit{rectified linear unit} (\textit{ReLU})
    \end{itemize}
\end{frame}

\begin{frame}{\textit{ResNet152V2}}
    A figura \ref{fig:residual-units}a representa a proposta original (representado na equação \ref{eq:residual-units}) e a figura \ref{fig:residual-units}b representa o que foi proposto para utilização nessa arquitetura. 
    \vspace{1em}
    
    Pode-se observar que a informação passa de uma forma mais fácil através da seta em cinza.

    \vspace{1em}
    Isso faz com que não possua problemas com a propagação da informação, ou seja, que a informação não seja desaparecida no final da rede \cite{he2016identitymappingsdeepresidual}.
\end{frame}

\begin{frame}[plain]
    \begin{figure}
        \centering
        \caption{Representação dos \textit{Residual units}}
        \includegraphics[trim={0 0 13.5cm 0},clip,height=7cm]{images/cnn-architecture/ResNet152V2/teaser.pdf}
        \begin{center}
            {\footnotesize Fonte: \textcite{he2016identitymappingsdeepresidual}}
        \end{center}
        \label{fig:residual-units}
    \end{figure}
\end{frame}

\begin{frame}{\textit{DensetNet201}}
    Abreviação para \textit{Dense Convolutional Network}, foi proposto por \textcite{huang2018denselyconnectedconvolutionalnetworks}.

    \vspace{1em}
    Um dos problemas que as \textit{CNN} muito profundas, ou seja, que possuem muitas camadas, enfrentam é a perda de informação quando chaga às camadas finais da \textit{CNN}.
\end{frame}

\begin{frame}[plain]
    \begin{figure}
        \centering
        \caption{Representação dos \textit{dense blocks}}
        \includegraphics[width=0.5\linewidth]{images/cnn-architecture/DenseNet201/denseBlock.pdf}
        \begin{center}
            Fonte: \textcite{huang2018denselyconnectedconvolutionalnetworks}
        \end{center}
        \label{fig:dense blocks}
    \end{figure}
\end{frame}

\begin{frame}{\textit{DensetNet201}}
    Como pode-se observar na figura \ref{fig:dense blocks}, toda a saída de um conjunto de camadas vai para a entrada do conjunto de camada subsequente, isso é denominado \textit{dense connectivity} por \textcite{huang2018denselyconnectedconvolutionalnetworks}.

    \vspace{1em}
    Assim a informação não há perda de informação para as camadas mais profundas da rede.
\end{frame}

\begin{frame}{\textit{Composite function}}
    Cada \textit{dense block} (formado por $n$ camadas) passa por uma camada de transição $\mathcal{H}_{l}(\cdot)$ (representado na figura \ref{fig:densenet}), a ``\textit{convolution}'' representa o \textit{composite function} que possui três operações sequenciais \cite{huang2018denselyconnectedconvolutionalnetworks}:

    \begin{itemize}
    \item \textit{batch normalization}
    \item \textit{rectified linear unit} (\textit{ReLU})
    \item Uma convolução de $3\times 3$
    \end{itemize}

\end{frame}


\begin{frame}[plain]
    \begin{figure}
        \centering
        \caption{Exemplo de uma \textit{DenseNet} com três \textit{dense blocks}, entre os blocos há uma camada de transição}
        \includegraphics[width=\linewidth]{images/cnn-architecture/DenseNet201/denseNet.pdf}
        {\centering \footnotesize Fonte: \textcite{huang2018denselyconnectedconvolutionalnetworks}}
        \label{fig:densenet}
    \end{figure}
\end{frame}

\begin{frame}{\textit{DensetNet201}}
    A tabela \ref{tab:densenet-imagenet} mostra as variações da arquitetura \textit{DenseNet}. A arquitetura que será utilizado é o \textit{DenseNet201}.
\end{frame}

\begin{frame}[plain]
\renewcommand{\arraystretch}{1.1}
\begin{table}
\centering
\caption{Variações da arquitetura de DenseNet}
\resizebox{10cm}{!}{
\begin{tabular}{c|c|c|l|l|l}
\hline
Camadas& Tamanho da saída& DenseNet-121& \multicolumn{1}{c|}{DenseNet-169}& \multicolumn{1}{c|}{DenseNet-201}& \multicolumn{1}{c}{DenseNet-264} \\ \hline
Convolution & \cross{112} & \multicolumn{4}{c}{\cross{7} convolução, stride 2} \\ \hline
Pooling & \cross{56} & \multicolumn{4}{c}{\cross{3} max pool, stride 2} \\ \hline
\begin{tabular}[c]{@{}c@{}}Dense Block\\ (1)\end{tabular} & \cross{56} & \multicolumn{1}{l|}{\conv{6}} & \conv{6} & \conv{6} & \conv{6} \\ \hline
\multirow{2}{*}{\begin{tabular}[c]{@{}c@{}}Transition Layer\\ (1)\end{tabular}} & \cross{56} & \multicolumn{4}{c}{\cross{1} convolução} \\ \cline{2-6}
 & \cross{28} & \multicolumn{4}{c}{\cross{2} pooling médio, stride 2} \\ \hline
\begin{tabular}[c]{@{}c@{}}Dense Block\\ (2)\end{tabular} & \cross{28} & \multicolumn{1}{l|}{\conv{12}} & \conv{12}& \conv{12} & \conv{12} \\ \hline
\multirow{2}{*}{\begin{tabular}[c]{@{}c@{}}Transition Layer\\ (2)\end{tabular}} & \cross{28} & \multicolumn{4}{c}{\cross{1} convolução} \\ \cline{2-6}
 & \cross{14} & \multicolumn{4}{c}{\cross{2} pooling médio, stride 2} \\ \hline
\begin{tabular}[c]{@{}c@{}}Dense Block\\ (3)\end{tabular} & \cross{14} & \multicolumn{1}{l|}{\conv{24}} & \conv{32} & \conv{48} & \conv{64} \\ \hline
\multirow{2}{*}{\begin{tabular}[c]{@{}c@{}}Transition Layer\\ (3)\end{tabular}} & \cross{14} & \multicolumn{4}{c}{\cross{1} convolução} \\ \cline{2-6}
 & \cross{7} & \multicolumn{4}{c}{\cross{2} pooling médio, stride 2} \\ \hline
\begin{tabular}[c]{@{}c@{}}Dense Block\\ (4)\end{tabular} & \cross{7} & \multicolumn{1}{l|}{\conv{16}} & \conv{32} & \conv{32} & \conv{48} \\ \hline
\multirow{2}{*}{\begin{tabular}[c]{@{}c@{}}Classification\\ Layer\end{tabular}} & \cross{1} & \multicolumn{4}{c}{\cross{7} pooling médio global} \\ \cline{2-6}
 & & \multicolumn{4}{c}{1000D fully-connected, softmax} \\ \hline
\end{tabular}
}
\label{tab:densenet-imagenet}
\begin{center}
    \footnotesize Fonte: \textcite[tradução pelo autor]{huang2018denselyconnectedconvolutionalnetworks}
\end{center}
\end{table}
\end{frame}

\begin{frame}[plain]
\renewcommand{\arraystretch}{1.1}
\begin{table}
\centering
% \caption{Variações da arquitetura de DenseNet}
\resizebox{14cm}{!}{
\begin{tabular}{c|c|c|l|l|l}
\hline
Camadas& Tamanho da saída& DenseNet-121& \multicolumn{1}{c|}{DenseNet-169}& \multicolumn{1}{c|}{DenseNet-201}& \multicolumn{1}{c}{DenseNet-264} \\ \hline
Convolution & \cross{112} & \multicolumn{4}{c}{\cross{7} convolução, stride 2} \\ \hline
Pooling & \cross{56} & \multicolumn{4}{c}{\cross{3} max pool, stride 2} \\ \hline
\begin{tabular}[c]{@{}c@{}}Dense Block\\ (1)\end{tabular} & \cross{56} & \multicolumn{1}{l|}{\conv{6}} & \conv{6} & \conv{6} & \conv{6} \\ \hline
\multirow{2}{*}{\begin{tabular}[c]{@{}c@{}}Transition Layer\\ (1)\end{tabular}} & \cross{56} & \multicolumn{4}{c}{\cross{1} convolução} \\ \cline{2-6}
 & \cross{28} & \multicolumn{4}{c}{\cross{2} pooling médio, stride 2} \\ \hline
\begin{tabular}[c]{@{}c@{}}Dense Block\\ (2)\end{tabular} & \cross{28} & \multicolumn{1}{l|}{\conv{12}} & \conv{12}& \conv{12} & \conv{12} \\ \hline
\multirow{2}{*}{\begin{tabular}[c]{@{}c@{}}Transition Layer\\ (2)\end{tabular}} & \cross{28} & \multicolumn{4}{c}{\cross{1} convolução} \\ \cline{2-6}
 & \cross{14} & \multicolumn{4}{c}{\cross{2} pooling médio, stride 2} \\ \hline
\begin{tabular}[c]{@{}c@{}}Dense Block\\ (3)\end{tabular} & \cross{14} & \multicolumn{1}{l|}{\conv{24}} & \conv{32} & \conv{48} & \conv{64} \\ \hline
\multirow{2}{*}{\begin{tabular}[c]{@{}c@{}}Transition Layer\\ (3)\end{tabular}} & \cross{14} & \multicolumn{4}{c}{\cross{1} convolução} \\ \cline{2-6}
 & \cross{7} & \multicolumn{4}{c}{\cross{2} pooling médio, stride 2} \\ \hline
\begin{tabular}[c]{@{}c@{}}Dense Block\\ (4)\end{tabular} & \cross{7} & \multicolumn{1}{l|}{\conv{16}} & \conv{32} & \conv{32} & \conv{48} \\ \hline
\multirow{2}{*}{\begin{tabular}[c]{@{}c@{}}Classification\\ Layer\end{tabular}} & \cross{1} & \multicolumn{4}{c}{\cross{7} pooling médio global} \\ \cline{2-6}
 & & \multicolumn{4}{c}{1000D fully-connected, softmax} \\ \hline
\end{tabular}
}
\end{table}
\end{frame}

\begin{frame}{\textit{InceptionV3}}
    Arquitetura de rede neural convolucional foi proposta por \textcite{szegedy2015rethinkinginceptionarchitecturecomputer}, possui um custo de processamento reduzido em comparação ao \textit{VGGNet} devido a sua menor quantidade de parâmetros para ser avaliado.
\end{frame}

\begin{frame}{\textit{InceptionV3}}
    Alguns princípios que foram adotados durante a concepção de sua arquitetura: evitar gargalos na representação das informações pelo processo de \textit{feed-forward} pela arquitetura da \textit{CNN}, ou seja, não fazer com que as informações sofram uma compressão extrema.
\end{frame}

\begin{frame}{\textit{InceptionV3}}
    Processamento em uma representação de alta dimensão é mais fácil para uma rede local e isso faz com que o treino seja mais rápido.
    
    \vspace{1em}
    Por último, balancear a \textit{CNN} no tamanho e profundidade \cite{szegedy2015rethinkinginceptionarchitecturecomputer}. 
\end{frame}

\begin{frame}{Utilização de convoluções menores de forma sequencial}
    Invés de utilizar uma convolução com o \textit{kernel} de tamanho de $5\times 5$ ou $7\times 7$, foi adotado a utilização de \textit{kernel} menor e utilizar sobre o resultado obtido na primeira passagem. 
    
    \vspace{1em}
    Isso apresenta uma redução no custo computacional para o treinamento e o processo de \textit{feed-forward} \cite{szegedy2015rethinkinginceptionarchitecturecomputer}.
\end{frame}



\begin{frame}[plain]
    \begin{figure}
        \centering
        \caption{Representação de duas convoluções de $3\times 3$}
        \includegraphics[width=0.5\linewidth]{images/cnn-architecture/InceptionV3/double3x3.pdf}
        \begin{center}
            Fonte: \textcite{szegedy2015rethinkinginceptionarchitecturecomputer}
        \end{center}
        \label{fig:enter-label}
    \end{figure}
\end{frame}

\begin{frame}{\textit{InceptionV3}}
    Esse método também apresenta um resultado equivalente a convoluções de $3\times 3$, porém utilizando, por exemplo, uma convolução de $1\times 3$ e em seguida $3\times 1$. 
    \vspace{1em}
    
    A diferença que utiliza 33\% a menos de processamento, caso o número de entrada e saída de filtros sejam iguais.
    \vspace{1em}
    
    Também propõem que pode-se utilizar convoluções de $1\times n$ seguido de $n\times 1$ \cite{szegedy2015rethinkinginceptionarchitecturecomputer}.
\end{frame}

\begin{frame}[plain]
    \begin{figure}
        \centering
        \caption{Representação de convoluções assimétricas}
        \includegraphics[width=0.5\linewidth]{images/cnn-architecture/InceptionV3/double31.pdf}
        \begin{center}
            Fonte: \textcite{szegedy2015rethinkinginceptionarchitecturecomputer}
        \end{center}
        \label{fig:enter-label}
    \end{figure}
    
\end{frame}



\begin{frame}{\textit{InceptionV3}}
    Para não remover características importantes e não ter um alto custo computacional, construíram um modulo que tem uma camada de \textit{pooling} e a outra de convolução e depois, o resultado dessas duas são concatenadas.
    \vspace{1em}

    Na figura \ref{fig:hybridreduction}, ambos exigem menos processamento e não removem características importantes \textcite{szegedy2015rethinkinginceptionarchitecturecomputer}.
\end{frame}

\begin{frame}[plain]
    \begin{figure}
        \centering
        \caption{Representação de uma alternativa para redução da dimensionalidade}
        \includegraphics[width=0.6\linewidth]{images/cnn-architecture/InceptionV3/hybridreduction.pdf}
        \begin{center}
            Fonte: \textcite{szegedy2015rethinkinginceptionarchitecturecomputer}
        \end{center}
        \label{fig:hybridreduction}
    \end{figure}
\end{frame}

\begin{frame}[plain]
    \begin{table}
        \caption{Diagrama da arquitetura do \textit{Inception}}
        {\small
         \begin{center}
            \resizebox{11cm}{!}{
           \begin{tabular}[H]{|l|c|c|}
           \hline
           {Tipo} & {Tamanho do \textit{kernel} ou \textit{strides}} & {Tamanho da entrada} \\
           \hline\hline
           conv & $3\times3/2$ & $299\times299\times3$ \\
           \hline
            conv & $3\times3/1$ & $149\times149\times32$ \\
           \hline
           conv padded & $3\times3/1$ & $147\times147\times32$ \\
           \hline
           pool & $3\times3/2$ & $147\times147\times64$ \\
           \hline
           conv & $3\times3/1$ & $73\times73\times64$ \\
           \hline
           conv & $3\times3/2$ & $71\times71\times80$ \\
           \hline
           conv & $3\times3/1$ & $35\times35\times192$ \\
           \hline
           $3\times$Inception & representado na figura~\ref{fig:inceptionv2} & $35\times35\times288$ \\
           \hline
           $5\times$Inception & representado na figura~\ref{fig:inceptionv3} & $17\times17\times768$ \\
           \hline
           $2\times$Inception & representado na figura~\ref{fig:inceptionv4} & $8\times8\times1280$ \\
           \hline
           pool & $8\times8$ & $8\times8\times2048$ \\
           \hline
           linear & logits & $1\times1\times2048$ \\
           \hline
           softmax& classifier & $1\times1\times1000$ \\
           \hline
           \end{tabular}}
         \end{center}
         }
         \begin{center}
            {\footnotesize Fonte: \textcite[tradução pelo autor]{szegedy2015rethinkinginceptionarchitecturecomputer}}
         \end{center}
        \label{tab:incetion}
        \end{table}
\end{frame}

\begin{frame}[plain]
    \begin{figure}
        \centering
        \caption{Módulos de \textit{Inception} utilizadas na arquitetura}
        \vspace{-0.5em}
        \begin{subfigure}[t]{0.3\textwidth}
             \centering
             \caption{Módulo \textit{Inception} com substituição de convolução de $5\times5$ por $3\times3$}
             \includegraphics[width=\textwidth]{images/cnn-architecture/InceptionV3/inceptionv2.pdf}
             \label{fig:inceptionv2}
         \end{subfigure}\hspace{\fill}
         \begin{subfigure}[t]{0.3\textwidth}
             \centering
             \caption{Módulos de \textit{Inception} para convoluções assimétricas}
             \includegraphics[width=\textwidth]{images/cnn-architecture/InceptionV3/inceptionv3.pdf}
             \label{fig:inceptionv3}
         \end{subfigure}\hspace{\fill}
         \begin{subfigure}[t]{0.3\textwidth}
             \centering
             \caption{Módulos de \textit{Inception} para expansão dos filtros de saída}
             \includegraphics[width=\textwidth]{images/cnn-architecture/InceptionV3/inceptionv4.pdf}
             \label{fig:inceptionv4}
         \end{subfigure}
         \vspace{-0.5em}
         \begin{center}
            {\footnotesize Fonte: \textcite{szegedy2015rethinkinginceptionarchitecturecomputer}}
         \end{center}
    \end{figure}
\end{frame}

\begin{frame}{\textit{InceptionResNetV2}}
    Arquitetura de rede neural convolucional foi proposta por \textcite{szegedy2016inceptionv4inceptionresnetimpactresidual}, na qual combina duas arquiteturas: 
    
    \begin{itemize}
        \item \textit{ResNet}, que foi proposta no artigo ``\textit{\citetitle{he2016identitymappingsdeepresidual}}''.
        \item \textit{Inception}, que foi proposta no artigo ``\textit{\citetitle{szegedy2015rethinkinginceptionarchitecturecomputer}}''.
    \end{itemize}    
\end{frame}
\begin{frame}{\textit{InceptionResNetV2}}
    Segundo \textcite{he2016identitymappingsdeepresidual}, a utilização de \textit{residual conections} demonstra uma vantagem para arquiteturas mais profundas.
    \vspace{1em}

    Para a concepção dessa arquitetura, foi utilizado os \textit{pure inception blocks} e \textit{residual inception blocks}.
\end{frame}
\begin{frame}{\textit{InceptionResNetV2}}
    A representação do diagrama da estrutura da arquitetura de \textit{CNN} do \textit{InceptionResNetV2} na figura \ref{fig:inceptionresnetv2}.
\end{frame}


\begin{frame}[plain]
    \begin{figure}
        \centering
        \caption{Diagrama da arquitetura de \textit{InceptionResNetV2}}
        \includegraphics[height=7cm]{images/cnn-architecture/InceptionResNetV2/inceptionresnetsmallschema.pdf}
        \label{fig:inceptionresnetv2}
        \begin{center}
            {\footnotesize Fonte: \textcite{szegedy2016inceptionv4inceptionresnetimpactresidual}}
        \end{center}
    \end{figure}
    
\end{frame}

\begin{frame}[plain]
    \begin{figure}
        \centering
        \begin{subfigure}[t]{0.45\textwidth}
            \centering
            \vspace{-55em}\includegraphics[width=\textwidth]{images/cnn-architecture/InceptionResNetV2/inceptionresnetsmallschema.pdf}
        \end{subfigure}
        \begin{subfigure}[t]{0.45\textwidth}
            \centering
            \includegraphics[width=\textwidth]{images/cnn-architecture/InceptionResNetV2/inceptionresnetsmallschema.pdf}
        \end{subfigure}
    \end{figure}
\end{frame}

\begin{frame}{\textit{InceptionResNetV2}}
    Onde:
    \begin{itemize}
        \item \textit{Stem}: representado na figura \ref{fig:inceptionv4stem}
        \item \textit{Inception-resnet-A}: representado na figura \ref{fig:resnetwide35x35}
        \item \textit{Reduction-A}: representado na figura \ref{fig:reductionto17}
        \item \textit{Inception-resnet-B}: representando na figura \ref{fig:resnetsmall17x17}
        \item \textit{Reduction-B}: representado na figura \ref{fig:reductionto8resnet}
        \item \textit{Inception-resnet-C}: representado na figura \ref{fig:resnetwide8x8}
    \end{itemize}
\end{frame}

\begin{frame}[plain]
    \begin{figure}
        \centering
            \caption{Representação do \textit{stem} do \textit{pure Inception}}
            \includegraphics[height=7cm]{images/cnn-architecture/InceptionResNetV2/inceptionv4stem.pdf}
            \label{fig:inceptionv4stem}
            \begin{center}
                {\footnotesize Fonte: \textcite{szegedy2016inceptionv4inceptionresnetimpactresidual}}
            \end{center}
    \end{figure}
\end{frame}

\begin{frame}[plain]
    \begin{figure}
        \centering
            \includegraphics[width=0.4\linewidth,trim={0 0 0 12cm},clip]{images/cnn-architecture/InceptionResNetV2/inceptionv4stem.pdf}
    \end{figure}
\end{frame}

\begin{frame}[plain]
    \begin{figure}
        \centering
            \includegraphics[width=0.4\linewidth,trim={0 7cm 0 0},clip]{images/cnn-architecture/InceptionResNetV2/inceptionv4stem.pdf}
    \end{figure}
\end{frame}

% \begin{frame}
%     \centering
%     \begin{figure}
%         \begin{subfigure}[t]{0.45\linewidth}
%             \centering
%             \includegraphics[width=\linewidth,trim={0 0 0 8cm},clip]{images/cnn-architecture/InceptionResNetV2/inceptionv4stem.pdf}
%         \end{subfigure}
%         \begin{subfigure}[t]{0.45\linewidth}
%             \centering
%             \includegraphics[width=\linewidth,trim={0 4cm 0 0},clip]{images/cnn-architecture/InceptionResNetV2/inceptionv4stem.pdf}
%         \end{subfigure}
%     \end{figure}
% \end{frame}

\begin{frame}[plain]
    \begin{figure}
        \centering
            \caption{Representação da convolução de $35\times 35$}
            \includegraphics[height=7cm]{images/cnn-architecture/InceptionResNetV2/resnetwide35x35.pdf}    
            \label{fig:resnetwide35x35}
            \begin{center}
                {\footnotesize Fonte: \textcite{szegedy2016inceptionv4inceptionresnetimpactresidual}}
            \end{center}
    \end{figure}
\end{frame}

\begin{frame}[plain]
    \begin{figure}
        \centering
            \caption{Representação da redução de dimensionalidade de $35\times 35$ para $17\times 17$}
            \includegraphics[height=7cm]{images/cnn-architecture/InceptionResNetV2/reductionto17.pdf}
            \label{fig:reductionto17}
            \begin{center}
                {\footnotesize Fonte: \textcite{szegedy2016inceptionv4inceptionresnetimpactresidual}}
            \end{center}
    \end{figure}
\end{frame}


\begin{frame}[plain]
    \begin{figure}
        \centering
        \caption{Representação da convolução de $17\times 17$}
        \includegraphics[height=7cm]{images/cnn-architecture/InceptionResNetV2/resnetsmall17x17.pdf}
        \label{fig:resnetsmall17x17}
    \begin{center}
        {\footnotesize Fonte: \textcite{szegedy2016inceptionv4inceptionresnetimpactresidual}}
    \end{center}
    \end{figure}
\end{frame}

\begin{frame}[plain]
    \begin{figure}
        \centering
        \caption{Representação da redução de dimensionalidade de $17\times 17$ para $8\times 8$}
        \includegraphics[height=7cm]{images/cnn-architecture/InceptionResNetV2/reductionto8resnet.pdf}
        \label{fig:reductionto8resnet}
        \begin{center}
            {\footnotesize Fonte: \textcite{szegedy2016inceptionv4inceptionresnetimpactresidual}}
            \end{center}
    \end{figure}
\end{frame}

\begin{frame}[plain]
    \begin{figure}
        \centering
            \caption{Representação da convolução de $8\times 8$}
            \includegraphics[height=7cm]{images/cnn-architecture/InceptionResNetV2/resnetwide8x8.pdf}
            \label{fig:resnetwide8x8}
            \begin{center}
            {\footnotesize Fonte: \textcite{szegedy2016inceptionv4inceptionresnetimpactresidual}}
            \end{center}
    \end{figure}
\end{frame}

\begin{frame}{\textit{MobileNetV2}}
    A arquitetura de \textit{CNN} foi proposta por \textcite{sandler2019mobilenetv2invertedresidualslinear}, o objetivo é conseguir utilizar a arquitetura de \textit{CNN} em dispositivos que possuem menos poder computacional (e.g. como celulares \textit{smartphones}, sistemas embarcados).
\end{frame}

\begin{frame}{\textit{MobileNetV2}}
    Alguns conceitos utilizados nessa arquitetura como \textit{depthwise convolution} foram introduzidos na primeira versão de \textit{MobileNet}, no artigo ``\textit{\citetitle{howard2017mobilenetsefficientconvolutionalneural}}''.
\end{frame}

\begin{frame}{\textit{MobileNetV2}}
    Conceitos como: \textit{linear bottlenecks} e \textit{inverted residuals} foram introduzidos na segunda versão da arquitetura (essa arquitetura).
\end{frame}

\begin{frame}{\textit{MobileNetV2}}
    O \textit{MobileNet} utiliza uma forma de convolução denominada \textit{depthwise separable convolutions} que utiliza de 8 a 9 vezes menos de processamento computacional.
    
    \vspace{1em}
    Essa camada de convolução é realizado em duas partes: \textit{depthwise convolution} e \textit{pointwiese convolution} \cite{howard2017mobilenetsefficientconvolutionalneural}.
\end{frame}

\begin{frame}{\textit{MobileNetV2}}
    O processo de \textit{Depthwise convolution} aplica uma convolução (representado na figura \ref{fig:depthwise separable convolutions}) que contém um único filtro com dimensão $D_{K}$ em cada canal da imagem \cite{howard2017mobilenetsefficientconvolutionalneural}.
\end{frame}

\begin{frame}[plain]
    \begin{figure}[H]
        \centering
        \caption{Representação de \textit{depthwise convolution}}
        \includegraphics[width=0.8\linewidth]{images/cnn-architecture/MobileNetV2/depthwise.pdf}
        \begin{center}
            {\footnotesize Fonte: \textcite{howard2017mobilenetsefficientconvolutionalneural}}
        \end{center}
        \label{fig:depthwise separable convolutions}
    \end{figure}
\end{frame}

\begin{frame}{\textit{MobileNetV2}}
    O processo de \textit{Pointwise convolution} aplica uma convolução (representado na figura \ref{fig:pointwise convolution}) com dimensão de $1\times 1$ com $M$ filtros combinando as saídas do \textit{depthwise convolution} \cite{howard2017mobilenetsefficientconvolutionalneural}.    
\end{frame}

\begin{frame}[plain]
    \begin{figure}
        \centering
        \caption{Representação de \textit{pointwise convolution}}
        \includegraphics[width=0.8\linewidth]{images/cnn-architecture/MobileNetV2/pointwise.pdf}
        \begin{center}
            {\footnotesize Fonte: \textcite{howard2017mobilenetsefficientconvolutionalneural}}
        \end{center}
        \label{fig:pointwise convolution}
    \end{figure}
\end{frame}

\begin{frame}{\textit{MobileNetV2}}
    Os \textit{bottlenecks blocks} estão representados na tabela \ref{tab:bottlenec_block_table} e na figura \ref{fig:linear bottlenecks} \cite{sandler2019mobilenetv2invertedresidualslinear}.
\end{frame}

\begin{frame}[plain]
    \begin{table}
        \centering
        \caption{Representação do \textit{bottleneck residual block}}
            \resizebox{11cm}{!}{
                \begin{tabular}{c|c|c}
                \hline
                Entrada & Operação & Saída\\
                \hline
                $h \times w \times k$ & $1\times 1$ \conv, \relus & $h \times w \times (tk)$\\
                $h \times w \times tk$& $3\times 3$ \depthwise s=$s$, \relus & $\frac{h}{s} \times \frac{w}{s} \times (tk)$\\ 
                $\frac{h}{s} \times \frac{w}{s} \times tk$ & linear $1\times 1$ \conv & $\frac{h}{s} \times \frac{w}{s} \times k'$\\
                \hline
                \end{tabular}}
            \begin{center}
                {\footnotesize Fonte: \textcite[tradução pelo autor]{sandler2019mobilenetv2invertedresidualslinear}}
            \end{center}
        \label{tab:bottlenec_block_table}
        \end{table}
\end{frame}

\begin{frame}[plain]
    \begin{figure}
        \centering
        \caption{Representação de \textit{linear bottleneck}}
        \includegraphics[width=\linewidth]{images/cnn-architecture/MobileNetV2/bottleneck_conv.pdf}
        \begin{center}
            {\footnotesize Fonte: \textcite{sandler2019mobilenetv2invertedresidualslinear}}
        \end{center}
        \label{fig:linear bottlenecks}
    \end{figure}
\end{frame}

\begin{frame}{\textit{MobileNetV2}}
    \textit{Inverted residuals} pode ser comparado ao \textit{residual units} (demonstrado no slide \ref{sec:ResNet152V2}) \cite{sandler2019mobilenetv2invertedresidualslinear}.
\end{frame}

\begin{frame}[plain]
    \begin{figure}
        \centering
        \caption{Representação de \textit{inverted residuals}}
        \includegraphics[width=\linewidth]{images/cnn-architecture/MobileNetV2/inverted_residual.pdf}
        \begin{center}
            {\footnotesize Fonte: \textcite{sandler2019mobilenetv2invertedresidualslinear}} 
        \end{center}
        \label{fig:inverted residuals}
    \end{figure}
\end{frame}

% \begin{frame}{\textit{MobileNetV2}}
%     Representação do diagrama da arquitetura \textit{MobileNetV2} na tabela \ref{tab:mobilenetarch}.
% \end{frame}

\begin{frame}[plain]
    \begin{table}
        \centering
        \caption{Diagrama da arquitetura \textit{MobileNetV2}}
        % \mbox{MobileNetV2} : Each line describes a sequence of 1 or more identical (modulo stride)  layers, repeated $n$ times.
        %     All layers in the same sequence have the same number $c$ of output channels.
        %     The first layer of each sequence has a stride $s$ and all others use stride $1$.
        %     All spatial convolutions use $3\times3$ kernels. The expansion factor $t$ is always applied to the input size as described in Table~\ref{fig:bottlenec_block_table}.
        \vspace{0pt}
        \begin{tabular}{c|c|c|c|c|c}
        % \toprule[0.2em]
        \hline
        Entrada & Operação & $t$& $c$ & $n$ & $s$ \\
        \hline
        % \toprule[0.2em]
            $224\times224\times3$ & conv2d & - & 32 & 1 & 2\\
            $112\times122\times32$ & bottleneck & 1 & 16 & 1 & 1 \\
            $112\times112\times16$ & bottleneck & 6 & 24 & 2 & 2 \\
            $56\times56\times24$ & bottleneck & 6 & 32 & 3 & 2 \\
            $28\times28\times32$ & bottleneck & 6 & 64 & 4 & 2 \\
            $14\times14\times64$ & bottleneck & 6 & 96 & 3 & 1 \\
            $14\times14\times96$ & bottleneck & 6 & 160 & 3 & 2 \\
            $7\times7\times160$ & bottleneck & 6 & 320 & 1 & 1 \\
            $7\times7\times320$ & conv2d 1$\times$1 & - & 1280 & 1 & 1 \\
            $7\times7\times1280$ & avgpool 7$\times$7 & - & - & 1 & - \\ 
            $1\times1\times1280$ & conv2d 1$\times$1 & - & k & -& \\
        % \toprule[0.2em]
        \hline
        \end{tabular}
            \begin{center}
                {\footnotesize Fonte: \textcite[tradução pelo autor]{sandler2019mobilenetv2invertedresidualslinear}}
            \end{center}
        \label{tab:mobilenetarch}
        \end{table}
\end{frame}

\begin{frame}{\textit{MobileNetV2}}
    Onde:
    \begin{itemize}
        \item c: representação do tamanho da saída 
        \item s: \textit{stride} da camada
        \item t: é o fator de expansão de cada saída
    \end{itemize}
\end{frame}